\documentclass{article}
\usepackage{geometry}
\usepackage[utf8]{inputenc}
\usepackage{hyperref}
\usepackage[table,xcdraw]{xcolor}
\usepackage{enumerate}
\usepackage{enumitem}
\usepackage{background}
\usepackage{graphicx}
\usepackage[a4paper, margin=0.9in]{geometry}
\usepackage{tikz}
\usetikzlibrary{calc}
\backgroundsetup{contents=\includegraphics{Geet-Group.jpg},opacity=0.7,angle=0,scale=1}
\geometry{left=1in, right=1in, top=1in, bottom=1in}
\usepackage{titling}

\begin{document}
\begin{tikzpicture}[remember picture,overlay]
    \draw[line width=4pt] 
        ($(current page.north west) + (0.5in,-0.5in)$) 
        rectangle 
        ($(current page.south east) + (-0.5in,0.5in)$);
\end{tikzpicture}

\begin{center}
    \huge\textbf{MAULANA ABUL KALAM AZAD \\ \vspace{0.5cm}UNIVERSITY OF TECHNOLOGY}\\
    \vspace{0.5 cm}{MAKAUT WB}

\end{center}
\begin{figure}[h!]
    \centering
    \includegraphics[width=0.3\linewidth]{MAKAUT_LOGO.jpg}
\end{figure}
\date{\today} 
\begin{center}
    \section*{\textbf{\underline{Software Tools and Technology}}}
    \vspace{0.5cm}
    \textbf{\underline{\large{Lab Notebook}}} \\
    \vspace{0.5cm}
    \LARGE submitted by: - \textbf{Group 25}
\end{center}
\begin{center}
\section*{\textbf{\underline{Group Members}}}
\end{center}
\vspace{0.2 cm}
\begin{tabbing}
    
    \hspace{5cm} \= \hspace{7cm} \= \kill
    \textbf{1. \underline{Subhankar Prasad Harijan}} (\textbf{LEADER}) \\
    \hspace{0.6cm}Department: - BCA \\
     \hspace{0.6 cm}Roll No: - 30001223005 \\
    
    \\
    \textbf{2. \underline{Gloria Ghosh}} \\
     \hspace{0.6 cm}Department: - BSc in IT (Data Science) \\
    \hspace{0.5 cm} Roll No: - 30084323024 \\
    \\
    \textbf{3. \underline{Sresta Tewari}} \\
    \hspace{0.4 cm} Department: - BSc in IT (Artificial Intelligence) \\
    \hspace{0.4 cm} Roll No: - 30054623004 \\
    \\
    \textbf{4. \underline{Sourav Sasmal}} \\
    \hspace{0.4 cm} Department: - BCA \\
    \hspace{0.4 cm} Roll No: - 30001223073 \\
    \\
    \textbf{5. \underline{Aakash Roy}} \\
     \hspace{0.6 cm}Department: - BSc in IT (Data Science) \\
    \hspace{0.5 cm} Roll No: - 30084323015 \\
\end{tabbing}


\newpage
\begin{tikzpicture}[remember picture,overlay]
    \draw[line width=3pt] 
        ($(current page.north west) + (0.5in,-0.5in)$) 
        rectangle 
        ($(current page.south east) + (-0.5in,0.5in)$);
\end{tikzpicture}


\vspace{3cm} 
\begin{center}
    \Huge \textbf{\textcolor{purple}{\underline{INDEX}}} 
\end{center}

\vspace{2cm} 

\begin{center}
\begin{tabular}{|>{\centering\arraybackslash}m{2cm}|m{8cm}|}
\hline
\textbf{Sl. No.} & \textbf{Questions} \\
\hline
1 & Introduction to GitHub and GitHub Desktop version installation. \\

\hline
2 & Building a C program for a calculator in the local repository, committing, and publishing it as a public repository. \\

\hline
3 & How to create matrix in LaTeX. \\

\hline
4 & How to create section, subsection, and subsubsection in LaTeX. \\

\hline
5 & Enumerate ABC Format and Roman Number in LaTeX.\\
\hline
\end{tabular}
\end{center}



\newpage 
\begin{tikzpicture}[remember picture,overlay]
    \draw[line width=4pt] 
        ($(current page.north west) + (0.5in,-0.5in)$) 
        rectangle 
        ($(current page.south east) + (-0.5in,0.5in)$);
\end{tikzpicture}

\begin{center}
    \Large{\textbf{\underline{Lab Assignment-1}}
    \end{center}
    \vspace{1cm}
   \begin{center}
    \textbf{\underline{ Introduction to GitHub and GitHub Desktop version installation}}
\end{center}
\begin{figure}[h!]
    \centering
    \includegraphics[width=0.3\linewidth]{The_geet.jpg}
\end{figure}
\begin{center}
    \section*{\textbf{\underline{GitHub}}}

\end{center}


\large{GitHub is a web-based platform that allows developers to host, share, and collaborate on software projects. It provides a version control system powered by Git, enabling teams to track changes, manage code repositories, and work together seamlessly across different locations. GitHub supports collaborative development through features like pull requests, issues, and project boards, making it essential for open-source projects and professional software development. Additionally, it integrates with various development tools, enhancing productivity and streamlining the software development lifecycle.}
\vspace{1 cm}
\begin{center}
    \section*{\textbf{\underline{Installation}}}
\end{center}

\large{Installing GitHub Desktop is a straightforward process that enhances your workflow by providing a user-friendly interface for managing repositories. To begin, download the installer from the \href{https://desktop.github.com/}{official GitHub Desktop website} for your operating system—Windows or macOS. After downloading, run the installer and follow the on-screen instructions to complete the setup. Once installed, launch the application and sign in with your GitHub credentials, or create a new account if needed. GitHub Desktop simplifies the process of cloning repositories, making commits, and managing branches, making it an invaluable tool for developers of all skill levels. For Linux users, alternative methods like using Wine or other Git clients are available.}

\newpage
\begin{tikzpicture}[remember picture,overlay]
    \draw[line width=3pt] 
        ($(current page.north west) + (0.5in,-0.5in)$) 
        rectangle 
        ($(current page.south east) + (-0.5in,0.5in)$);
\end{tikzpicture}

\begin{center}
    \Large{\textbf{\underline{Lab Assignment-2}}}
\end{center}
\section{\underline{Calculator Program using C}}
\subsection{\underline{Objective}}
\LARGE{The objective of this lab is to develop a basic calculator program using the C programming language. The calculator will perform simple arithmetic operations like addition, subtraction, multiplication, and division based on user input.}

\subsection{\underline{Program Overview}}
The calculator program is designed to:
\begin{itemize}
    \item Accept two numbers from the user.
    \item Prompt the user to select an arithmetic operation (Addition, Subtraction, Multiplication, Division).
    \item Perform the selected operation.
    \item Display the result of the operation to the user.
\end{itemize}

The program includes error handling to manage division by zero and other invalid inputs.

\subsection{\underline{Code Implementation}}
The following is the C code for the calculator program:

\newpage
\begin{tikzpicture}[remember picture,overlay]
    \draw[line width=3pt] 
        ($(current page.north west) + (0.5in,-0.5in)$) 
        rectangle 
        ($(current page.south east) + (-0.5in,0.5in)$);
\end{tikzpicture}

\begin{verbatim}
#include<stdio.h>
#include<conio.h>
void main(){
    float a,b,c;
    char ch;
    printf("Enter the first number : ");
    scanf("%f",&a);
    printf("Enter user choice operation : ");
    scanf(" %c",&ch);
    printf("Enter the second number : ");
    scanf("%f",&b);
    switch(ch){
        case'+':c=a+b;
        printf("Output is : %f",c);
        break;
        case'-':c=a-b;
        printf("Output is : %f",c);
        break;
        case'*':c=a*b;
        printf("Output is : %f",c);
        break;
        case'/':c=a/b;
        printf("Output is : %f",c);
        break;
        default:printf("Invalid operation");
        break;
    }
    getch();
}
\end{verbatim}
\begin{tikzpicture}[remember picture,overlay]
    \draw[line width=3pt] 
        ($(current page.north west) + (0.5in,-0.5in)$) 
        rectangle 
        ($(current page.south east) + (-0.5in,0.5in)$);
\end{tikzpicture}

\subsection{\underline{Compiling and Running the Program}}
To compile and run the calculator program:
\begin{enumerate}
    \item Open a terminal or command prompt.
    \item Navigate to the directory where the C file is located.
    \item Compile the program using a C compiler (e.g., GCC):
    \begin{verbatim}
    gcc calculator.c -o calculator
    \end{verbatim}
    \item Run the compiled program:
    \begin{verbatim}
    ./calculator
    \end{verbatim}
\end{enumerate}

\subsection{\underline{Adding the Calculator Program to GitHub Repository}}
To add this calculator program to a GitHub repository, follow these steps:

\subsubsection{Step 1: Initialize a Local Git Repository}
\begin{enumerate}
    \item Open the terminal and navigate to the directory where your \texttt{calculator.c} file is located.
    \item If you haven't already, initialize a Git repository in that directory:
    \begin{verbatim}
    git init
    \end{verbatim}
    This command creates a new Git repository in the current directory.
\end{enumerate}
\newpage
\begin{tikzpicture}[remember picture,overlay]
    \draw[line width=3pt] 
        ($(current page.north west) + (0.5in,-0.5in)$) 
        rectangle 
        ($(current page.south east) + (-0.5in,0.5in)$);
\end{tikzpicture}

\subsubsection{Step 2: Add the File to the Repository}
\begin{enumerate}
    \item Add the \texttt{calculator.c} file to the staging area:
    \begin{verbatim}
    git add calculator.c
    \end{verbatim}
    This command stages the file, indicating that you want to include it in the next commit.
\end{enumerate}

\subsubsection{Step 3: Commit the Changes}
\begin{enumerate}
    \item Commit the file to the repository with a meaningful message:
    \begin{verbatim}
    git commit -m "Add calculator program in C"
    \end{verbatim}
\end{enumerate}

\begin{tikzpicture}[remember picture,overlay]
    \draw[line width=3pt] 
        ($(current page.north west) + (0.5in,-0.5in)$) 
        rectangle 
        ($(current page.south east) + (-0.5in,0.5in)$);
\end{tikzpicture}

\subsubsection{Step 4: Push the Changes to GitHub}
\begin{enumerate}
    \item Link your local repository to a remote GitHub repository:
    \small{
    \begin{verbatim}
    git remote add origin https://github.com/yourusername/your-repo-name.git
    \end{verbatim} }
    \item Push the changes to the GitHub repository:
    \begin{verbatim}
    git push -u origin master
    \end{verbatim}
\end{enumerate}

\subsubsection{Step 5: Verify the Upload}
\begin{enumerate}
    \item Go to your GitHub repository URL in a web browser.
    \item Verify that the \texttt{calculator.c} file is listed and accessible in the repository.
\end{enumerate}


\begin{tikzpicture}[remember picture,overlay]
    \draw[line width=4pt] 
        ($(current page.north west) + (0.5in,-0.5in)$) 
        rectangle 
        ($(current page.south east) + (-0.5in,0.5in)$);
\end{tikzpicture}


\newpage
\begin{tikzpicture}[remember picture,overlay]
    \draw[line width=3pt] 
        ($(current page.north west) + (0.5in,-0.5in)$) 
        rectangle 
        ($(current page.south east) + (-0.5in,0.5in)$);
\end{tikzpicture}

\newpage
\begin{center}
    \LARGE{\textbf{\underline{Lab Assignment-3}}}
\end{center}
\vspace{0.2cm}
 \begin{center}
     \section*{\underline{How to create Matrix in LaTeX}}
\end{center}
\LARGE{To create a matrix in LaTeX, we use the \texttt{\textbf{amsmath}} package, which provides various environments for displaying matrices. Here’s a basic guide on how to create a matrix and the different options available:}
\vspace{0.1cm}
\begin{center}
\begin{itemize}
    \item{\underline{\textbf{Include the \texttt{amsmath} Package}}}
\end{itemize} 
\end{center}
First, ensure you have the \texttt{amsmath} package included in your LaTeX document preamble. Add the following line:
\begin{verbatim}
\usepackage{amsmath}
\end{verbatim}
\begin{center}
    \begin{itemize}
        \item {\textbf{\underline{Matrix Environments}}}
    \end{itemize}
\end{center}
\LARGE{
There are several environments for creating matrices, depending on how you want them formatted. Here are the most common ones:}

\begin{itemize}
  \item \texttt{matrix}: A basic matrix without brackets.
  \item \texttt{bmatrix}: A matrix with square brackets.
  \item \texttt{pmatrix}: A matrix with parentheses.
  \item \texttt{vmatrix}: A matrix with vertical bars.
  \item \texttt{Vmatrix}: A matrix with double vertical bars.
\end{itemize}
\begin{center}
        \item {\textbf{\underline{\LARGE{Syntax}}}}
\end{center}

\LARGE{The syntax for creating a matrix is similar across different environments. Use the \texttt{\textbackslash begin\{environment\}} and \texttt{\textbackslash end\{environment\}} commands to enclose the matrix content. Separate the elements in each row with \texttt{\&} and 

\begin{tikzpicture}[remember picture,overlay]
    \draw[line width=4pt] 
        ($(current page.north west) + (0.5in,-0.5in)$) 
        rectangle 
        ($(current page.south east) + (-0.5in,0.5in)$);
\end{tikzpicture}
end each row with \texttt{\\}.

Here’s an example of a 2x2 matrix in each environment:}
\begin{tikzpicture}[remember picture,overlay]
    \draw[line width=4pt] 
        ($(current page.north west) + (0.5in,-0.5in)$) 
        rectangle 
        ($(current page.south east) + (-0.5in,0.5in)$);
\end{tikzpicture}
\begin{verbatim}
[ 
$ \begin{matrix}
    1 & 0 \\
    0 & 1 
\end{matrix}
$ ]

\begin{bmatrix}   [
 $   \begin{matrix}
    1 & 0\\
    0 & 1
    \end{matrix}
 $   ]

a_{11} & a_{12} \\
a_{21} & a_{22}
\end{bmatrix}

\begin{pmatrix}
a_{11} & a_{12} \\
a_{21} & a_{22}
\end{pmatrix}

\begin{vmatrix}
a_{11} & a_{12} \\
a_{21} & a_{22}
\end{vmatrix}

\begin{Vmatrix}
a_{11} & a_{12} \\
a_{21} & a_{22}
\end{Vmatrix}
\end{verbatim}
\begin{tikzpicture}[remember picture,overlay]
    \draw[line width=3pt] 
        ($(current page.north west) + (0.5in,-0.5in)$) 
        rectangle 
        ($(current page.south east) + (-0.5in,0.5in)$);
\end{tikzpicture}

\begin{center}
\section*{\textbf{\underline{\LARGE{Explanation}}}}
\end{center}
\begin{itemize}
  \item \texttt{\textbackslash begin\{matrix\} ... \textbackslash end\{matrix\}}: Creates a matrix with no surrounding brackets.
  \item \texttt{\textbackslash begin\{bmatrix\} ... \textbackslash end\{bmatrix\}}: Surrounds the matrix with square brackets.
  \item \texttt{\textbackslash begin\{pmatrix\} ... \textbackslash end\{pmatrix\}}: Surrounds the matrix with parentheses.
  \item \texttt{\textbackslash begin\{vmatrix\} ... \textbackslash end\{vmatrix\}}: Surrounds the matrix with single vertical bars, often used to denote determinants.
  \item \texttt{\textbackslash begin\{Vmatrix\} ... \textbackslash end\{Vmatrix\}}: Surrounds the matrix with double vertical bars, also used for determinants or norms.
\end{itemize}

We can adjust the number of rows and columns by adding or removing elements and rows as needed.



\end{document}









